% Generated draft: requires mathematical and source review.
\documentclass[11pt]{article}
\usepackage[margin=25mm]{geometry}
\usepackage{fontspec}
\setmainfont{Latin Modern Roman}
\usepackage{amsmath,amssymb,mathtools,unicode-math}
\setmathfont{Latin Modern Math}
\usepackage{enumitem,needspace,xurl,hyperref}
\setlength{\parindent}{0pt}
\setlength{\parskip}{5pt}
\setlength{\emergencystretch}{3em}
\newcommand{\N}{\mathbb{N}}
\newcommand{\Z}{\mathbb{Z}}
\newcommand{\Q}{\mathbb{Q}}
\newcommand{\R}{\mathbb{R}}
\newcommand{\C}{\mathbb{C}}
\newcommand{\F}{\mathbb{F}}
\newcommand{\A}{\mathbb{A}}
\newcommand{\PP}{\mathbb{P}}
\newcommand{\E}{\mathbb{E}}
\newcommand{\eps}{\varepsilon}
\newcommand{\dd}{\,\mathrm{d}}
\begin{document}
\section*{UnsolvedMath ID: 30003731}
\textbf{Source problem code:} catalogue entry
\subsection{Definitions and mathematical statement}
The supplied record asks about Monte Carlo adjusted profile confidence intervals constructed from likelihood estimates obtained by iterated filtering and sequential Monte Carlo, especially for the model classes called PanelPOMPs and GenPOMPs. It does not define these classes, the interval construction, or an asymptotic regime. Consequently, the following expresses the intended coverage question schematically rather than specifying a complete conjecture.

For a coverage formulation, let \(\{P_{n,\theta}:\theta\in\Theta_n\}\) denote a sequence of statistical experiments, where \(n\) is an unspecified asymptotic index. Suppose a real-valued target \(\psi_n:\Theta_n\to\R\) has been selected. Let \(Y_n\) denote the observed data, and let \(U_n\) represent all randomness used in the Monte Carlo calculations. For a nominal error level \(\alpha\in(0,1)\), write
\[
C_{n,\alpha}(Y_n,U_n)
\]
for the interval returned by the intended Monte Carlo adjusted profile procedure using likelihood estimates obtained by iterated filtering and sequential Monte Carlo. These symbols introduce placeholders; neither the target nor the procedure is specified by the source.

A usual interpretation of asymptotically valid coverage would require, for every fixed \(\alpha\in(0,1)\) and every admissible sequence of true parameters \((\theta_n)\),
\[
\liminf_{n\to\infty}
\Pr_{n,\theta_n}\!\left\{
\psi_n(\theta_n)\in C_{n,\alpha}(Y_n,U_n)
\right\}
\geq 1-\alpha,
\]
where \(\Pr_{n,\theta_n}\) includes both sampling randomness and computational randomness. The admissible parameter sequences, the computational budgets, and their dependence on \(n\) must be specified before this becomes a definite assertion. The source does not determine whether coverage should instead be conditional on computational randomness, pointwise in a fixed parameter, uniform over a parameter class, or asymptotically equal to \(1-\alpha\).

\textbf{Source question.} Under what precisely stated model assumptions, regularity conditions, asymptotic regime, and Monte Carlo budget requirements do the intended intervals have asymptotically valid coverage for large models, including the intended PanelPOMP and GenPOMP classes?

The original wording asks for ``asymptotically valid theoretical guarantees,'' while the supplied cleaned formulation specifically asks for coverage. The displayed coverage condition is therefore an interpretation of the cleaned question, not a claim that the broader original wording has a uniquely determined mathematical meaning. The supplied dated assessment describes asymptotic justification in restricted quadratic-profile regimes and treats general guarantees for large PanelPOMP and GenPOMP models as a conjectural extension; those restrictions cannot be removed on the basis of this record.
\subsection{Short English statement}
The question asks when Monte Carlo adjusted profile intervals, computed using iterated filtering and sequential Monte Carlo likelihood estimates, attain their nominal coverage in large PanelPOMP and GenPOMP models. A precise conjecture still requires definitions of the model classes, target parameter, interval algorithm, limiting regime, and computational budgets. The supplied curation labels the problem partially solved and reports justification in restricted quadratic-profile regimes; this is a dated source assessment, not an independently verified conclusion.
\subsection{Sources}
\begin{itemize}
\item Supplied local UnsolvedMath registry record, permanent ID \(30003731\), source problem code \texttt{catalogue entry}, title \emph{Adjusted Profile Confidence Intervals for Large Monte Carlo Models}, in the Oberwolfach Reports Open Problems collection. The record gives proposed year \(2018\) and names no proposer.
\item Supplied report citation: \emph{Design and Analysis of Infectious Disease Studies (2019)}, Oberwolfach Reports, DOI \texttt{10.4171/owr/2018/7}, pp.\ 383--432. Supplied source link: \url{https://doi.org/10.4171/owr/2018/7}. The background attributes the problem to workshop open-problem discussions documented in Oberwolfach Reports.
\item Supplied literature source labeled ``original OWR report'': \url{https://doi.org/10.4171/OWR/2018/7}.
\item Supplied literature source: Ionides et al., \emph{Monte Carlo profile confidence intervals}, \url{https://doi.org/10.1098/rsif.2017.0126}. Full author details and theorem statements are not supplied.
\item Literature assessment checked at \texttt{2026-08-21}: the supplied curation labels the problem partially solved, reports published methodology and asymptotic justification in restricted quadratic-profile regimes, and describes general guarantees for large PanelPOMP and GenPOMP models as a conjectural extension. This is supplied dated curation derived from the OWR record, not independently verified literature assessment. The source itself requests independent verification of time-sensitive status and cited claims.
\end{itemize}
\end{document}
